\noindent\textbf{Step 2.} There is a coefficient issue in the parabolic Inoue case
which the torsion hypothesis resolves explicitly. Write
$E$ for its elliptic curve and $C$ for its reduced rational
cycle. The exact canonical relation is
$K_{Z_0}+E+C=0$; the curves $E$ and $C$ are disjoint,
$E^2=-b_2(Z_0)$, and the intersection matrix of the
components of $C$ has kernel generated by their sum
\cite[Lemma~2.22 and the following paragraph; Proposition~2.27]{Dloussky2006}.
Consequently an effective integral divisor with
$c_1(K_{Z_0}+D_0)_{\mathbb R}=0$ has the form
$D_0=E+kC$, $k\ge0$: intersection with $E$ fixes its
coefficient at one, and the cycle intersection matrix
makes all cycle coefficients equal. Our stronger
torsion relation then gives
$K_{Z_0}+D_0=(k-1)C$ torsion. For any Gauduchon metric
$\omega$, the degree of $\OO(C)$ is
$\int_C\omega>0$, whereas a torsion line has degree
zero. Hence $k=1$. In particular both connected
pieces of the stated reduced support really occur.
This deduction uses the torsion relation, not merely
the vanishing of its real first Chern class.

